\documentclass[10pt,a4paper]{amsart}
\usepackage[margin=19mm]{geometry}
\usepackage{amsmath,amssymb,mathtools}
\usepackage[scr=rsfs,cal=euler]{mathalfa}
\usepackage{xfrac}
\usepackage[unicode,hidelinks,bookmarks=false]{hyperref}
\newcommand{\ie}{i.e.\ }
\newcommand{\cf}{cf.\ }
\newcommand{\resp}{respectively\ }
\newcommand{\Subsection}{\subsection}
\providecommand{\nobreakdash}{\nobreak\hbox{-}}
\setlength{\emergencystretch}{2em}
\multlinegap0pt
\usepackage{amssymb}
\usepackage{xcolor}
\newtheorem{lemma}{Lemma}
\newtheorem{corollary}{Corollary}
\newtheorem{theorem}{Theorem}
\title{New upper bounds for the period of a negative orientable sequence}
\author{Chris J Mitchell, Peter R Wild}
\date{Source: arXiv:2602.04433v2 (CC BY 4.0)}
\begin{document}
\maketitle

\noindent\emph{Source and licence.} New upper bounds for the period of a negative orientable sequence. Version arXiv:2602.04433v2 (CC BY 4.0), distributed by arXiv under the Creative Commons Attribution 4.0 International licence, http://creativecommons.org/licenses/by/4.0/, as recorded in the arXiv licence metadata for this exact version and shown on the abstract page https://arxiv.org/abs/2602.04433v2. Full licence notice: \texttt{licenses/arxiv-2602.04433-CC-BY-4.0.txt}.\par\smallskip

\noindent\textit{Editorial scope.} Counted unit: the article's theorem theorem:new\_NOS\_bounds (source0.tex lines 707--733). The article's complete original proof of it is delivered with it (source0.tex lines 735--874, one \texttt{proof} environment of 140 lines). No mathematics is written, repaired or re-derived: the statement and the proof are the article's own lines, in the article's own notation, and no retained line is substituted or re-flowed.  Retained uncounted as the source-local context the proof needs: lines 203--250; lines 299--306; lines 423--445; lines 521--534; lines 550--580; lines 670--687.  Nothing was omitted from any retained span, so the package carries no omission record.  No retained line contains a citation, so the wrapper carries no bibliography and no bibliography entry.  No retained line contains a figure environment, an image-inclusion command or drawing code, so no image is delivered and the package declares no asset dependency.  Editorial formatting only, in addition: an AMS \texttt{amsart} preamble that restates the article's own declarations and macros used by the retained lines, namely the environments \texttt{lemma}, \texttt{corollary}, \texttt{theorem} on separate counters, together with the standard packages those lines need.  The retained environments print in this document as Lemma 1, Corollary 1, Theorem 1 in order of appearance, whereas the article prints the same items with the numbering of its own full document.  The complete native source of the article as distributed in the arXiv e-print for this exact version is bundled unchanged as \texttt{sources/193975/source0.tex}.
\begin{lemma}  \label{lemma:NOS_tuple_numbers}
Suppose $n\geq1$ and $k>2$.  Then:
\begin{itemize}
\item[i)]  the number of $k$-ary negasymmetric $n$-tuples is
\begin{align*}
k^{(n-1)/2} & \mbox{~~if $n$ is odd and $k$ is odd} \\
2k^{(n-1)/2} & \mbox{~~if $n$ is odd and $k$ is even;} \\
k^{n/2} & \mbox{~~if $n$ is even;}
\end{align*}
\item[ii)] the number of uniform $k$-ary $n$-tuples is $k$;
\item[iii)] the number of uniform-alternating $k$-ary $n$-tuples is $k$;
\item[iv)] the number of $n$-tuples that are both uniform and uniform-alternating is 1 if
    $k$ is odd and 2 if $k$ is even;
\item[v)] the number of uniform negasymmetric $k$-ary $n$-tuples is 1 if $k$ is odd and 2 if
    $k$ is even;
\item[vi)] the number of uniform-alternating negasymmetric $k$-ary $n$-tuples is 1 if $n$ and
    $k$ are odd, 2 if $n$ is odd and $k$ is even, and $k$ if $n$ is even;
\item[vii)] the number of alternating negasymmetric $k$-ary $n$-tuples is zero if $n$ and $k$
    are both odd, 2 if $n$ is odd and $k$ is even, $k-1$ if $n$ is even and $k$ is odd, and
    $k-2$ if $n$ and $k$ are both even (and in every case they are uniform-alternating);
\item[viii)] the number of left-sns $k$-ary $n$-tuples is:
\begin{align*}
k^{(n+1)/2} & \mbox{~~if $n$ is odd;} \\
k^{n/2} & \mbox{~~if $n$ is even and $k$ is odd;} \\
2k^{n/2} & \mbox{~~if $n$ and $k$ are both even;}
\end{align*}
\item[ix)] the number of non-uniform left-sns $k$-ary $n$-tuples is:
\begin{align*}
k^{(n+1)/2}-1 & \mbox{~~if $n$ and $k$ are both odd;} \\
k^{(n+1)/2}-2 & \mbox{~~if $n$ is odd and $k$ is even;} \\
k^{n/2}-1 & \mbox{~~if $n$ is even and $k$ is odd;} \\
2k^{n/2}-2 & \mbox{~~if $n$ and $k$ are both even;}
\end{align*}
\item[x)] the number of non-uniform-alternating left-sns $k$-ary $n$-tuples is:
\begin{align*}
k^{(n+1)/2}-k & \mbox{~~if $n$ is odd;} \\
k^{n/2}-1 & \mbox{~~if $n$ is even and $k$ is odd;} \\
2k^{n/2}-2 & \mbox{~~if $n$ and $k$ are both even.}
\end{align*}
\item[xi)] the number of non-uniform non-alternating left-sns $k$-ary $n$-tuples is:
\begin{align*}
k^{(n+1)/2}-k & \mbox{~~if $n$ is odd;} \\
k^{n/2}-1 & \mbox{~~if $n$ is even and $k$ is odd;} \\
2k^{n/2}-4 & \mbox{~~if $n$ and $k$ are both even;}
\end{align*}
\end{itemize}
Finally observe that (viii), (ix), (x) and (xi) also hold if left-sns is replaced with right-sns.
\end{lemma}

\begin{lemma}  \label{lemma:number_non-negasymmetric_tuples}
If $k>2$ and $n\geq2$ the number $N_k(n)$ of edges in $B_k^-(n-1)$ is
\begin{align*}
k^n-k^{(n-1)/2} & \mbox{~~if $n$ is odd and $k$ is odd} \\
k^n-2k^{(n-1)/2} & \mbox{~~if $n$ is odd and $k$ is even} \\
k^n-k^{n/2} & \mbox{~~if $n$ is even.}
\end{align*}
\end{lemma}

\begin{corollary}  \label{corollary:NOS_excluded_tuples_by_unequal_degree}
Suppose $k>2$ and $n\geq4$ and $S$ is an $\mathcal{NOS}_k(n)$. Then
\begin{itemize}
\item[i)] if $n$ and $k$ are both odd, for every vertex in $B^-_k(n-1)$ corresponding to a
    non-uniform left-sns $(n-1)$-tuple $(a_0,a_1,\ldots,a_{n-2})$, there is an edge
    $(a_0,a_1,\ldots,a_{n-2},x)$ in $B^-_k(n-1)$, for some $x$, that is not in $B^-(S,n)$.
\item[ii)] if $n$ and $k$ are both odd, for every vertex in $B^-_k(n-1)$ corresponding to a
    non-uniform right-sns $(n-1)$-tuple $(a_0,a_1,\ldots,a_{n-2})$, there is an edge
    $(y,a_0,a_1,\ldots,a_{n-2})$ in $B^-_k(n-1)$, for some $y$, that is not in $B^-(S,n)$.
\item[iii)] if $n$ is odd and $k$ is even, for every vertex in $B^-_k(n-1)$ corresponding to a
    non-uniform non-alternating left-sns $(n-1)$-tuple $(a_0,a_1,\ldots,a_{n-2})$, there is an
    edge $(a_0,a_1,\ldots,a_{n-2},x)$ in $B^-_k(n-1)$, for some $x$, that is not in $B^-(S,n)$.
\item[iv)] if $n$ is odd and $k$ is even, for every vertex in $B^-_k(n-1)$ corresponding to a
    non-uniform non-alternating right-sns $(n-1)$-tuple $(a_0,a_1,\ldots,a_{n-2})$, there is an
    edge $(y,a_0,a_1,\ldots,a_{n-2})$ in $B^-_k(n-1)$, for some $x$, that is not in $B^-(S,n)$.
\item[v)] if $n$ is even, for every vertex in $B^-_k(n-1)$ corresponding to a
    non-uniform-alternating left-sns $(n-1)$-tuple $(a_0,a_1,\ldots,a_{n-2})$, there is an edge
    $(a_0,a_1,\ldots,a_{n-2},x)$ in $B^-_k(n-1)$, for some $x$, that is not in $B^-(S,n)$.
\item[vi)] if $n$ is even, for every vertex in $B^-_k(n-1)$ corresponding to a
    non-uniform-alternating right-sns $(n-1)$-tuple $(a_0,a_1,\ldots,a_{n-2})$, there is an
    edge $(y,a_0,a_1,\ldots,a_{n-2})$ in $B^-_k(n-1)$, for some $y$, that is not in $B^-(S,n)$.
\end{itemize}
\end{corollary}

\begin{corollary}  \label{corollary:excluded-tuples-by-parity}
Suppose $k>2$ and $n\geq2$ and $S$ is an $\mathcal{NOS}_k(n)$. Then in all the following cases
there are edges $(a_0,a_1,\ldots,a_{n-2},x)$ and $(y,a_0,a_1,\ldots,a_{n-2})$ in $B^-_k(n-1)$
that are not in $B^-(S,n)$.
\begin{itemize}
\item[i)] if $k$ is odd, for every vertex in $B^-_k(n-1)$ corresponding to a negasymmetric
    non-uniform $(n-1)$-tuple $(a_0,a_1,\ldots,a_{n-2})$;
\item[ii)] if $n$ is odd and $k$ is even, for every vertex in $B^-_k(n-1)$ corresponding to a
    uniform or alternating negasymmetric $(n-1)$-tuple $(a_0,a_1,\ldots,a_{n-2})$, where
    $a_i\in\{0,k/2\}$ for every $i$;
\item[iii)] if $n$ and $k$ are both even, for every vertex in $B^-_k(n-1)$ corresponding to a
    uniform-alternating negasymmetric $(n-1)$-tuple $(a,-a,a,-a,\dots,a)$.
\end{itemize}
\end{corollary}

\begin{lemma}  \label{lemma:NOS-constraint-interactions}
Suppose $k>2$.  Suppose $\mathbf{a}=(a_0,a_1,\dots,a_{n-2})$ and
$\mathbf{b}=(a_1,a_2,\dots,a_{n-1})$ are $k$-ary $(n-1)$-tuples.
\begin{itemize}
\item[i)] If $n\geq3$, and $\mathbf{a}$ and $\mathbf{b}$ are both negasymmetric then
    \begin{itemize}
    \item if $n$ is even then $\mathbf{a}$ and $\mathbf{b}$ are both uniform or
        alternating, where $a_i\in\{0,k/2\}$ for every $i$;
    \item if $n$ is odd then $\mathbf{a}$ and $\mathbf{b}$ are uniform-alternating.
    \end{itemize}

\item[ii)] If $n\geq5$, $\mathbf{a}$ is negasymmetric and $\mathbf{b}$ is right-sns then there
    exist $c_j$, $0\leq j\leq2$, such that $a_{3i+j}=c_j$, for every $i$ and $0\leq j\leq 2$.
    Moreover if $n\equiv0\pmod3$ then $c_2\in\{0,k/2\}$ and $c_0=-c_1$; if $n\equiv1\pmod3$
    then $c_1\in\{0,k/2\}$ and $c_0=-c_2$; and if $n\equiv2\pmod3$ then $c_0\in\{0,k/2\}$ and
    $c_1=-c_2$.

\item[iii)] If $n\geq5$, $\mathbf{a}$ is left-sns and $\mathbf{b}$ is negasymmetric then there
    exist $c_j$, $0\leq j\leq2$, such that $a_{3i+j}=c_j$, for every $i$ and $0\leq j\leq 2$.
    Moreover if $n\equiv0\pmod3$ then $c_0\in\{0,k/2\}$ and $c_1=-c_2$; if $n\equiv1\pmod3$
    then $c_2\in\{0,k/2\}$ and $c_0=-c_1$; and if $n\equiv2\pmod3$ then $c_1\in\{0,k/2\}$ and
    $c_0=-c_2$.

\item[iv)] If $n\geq5$, $\mathbf{a}$ is left-sns and $\mathbf{b}$ is right-sns then there exist
    $c_j$, $0\leq j\leq3$, such that $a_{4i+j}=c_j$, for every $i$ and $0\leq j\leq 3$.
    Moreover if $n\equiv0\pmod4$ then $c_0=-c_1$ and $c_2=-c_3$; if $n\equiv1\pmod4$ then
    $c_0=-c_2$, $c_1\in\{0,k/2\}$, and $c_3\in\{0,k/2\}$; if $n\equiv2\pmod4$ then $c_0=-c_3$
    and $c_1=-c_2$; and if $n\equiv3\pmod4$ then $c_1=-c_3$, $c_0\in\{0,k/2\}$, and
    $c_2\in\{0,k/2\}$.
\end{itemize}
\end{lemma}

\begin{lemma}  \label{lemma:NOS-bounds-high-level}
If $k>2$ and $S$ is an $\mathcal{NOS}_k(n)$ then:
\[ |B^-(S,n)|\leq \begin{cases}
N_k(2) - |P_{\text{out}}|,
 & \text{if $n=2$} \\
N_k(3) - |P_{\text{out}}|-|P_{\text{in}}| + |P_{\text{out}}\cap P_{\text{in}}|,
 & \text{if $n=3$} \\
N_k(4) - |U_{\text{out}}|-|P_{\text{out}}|,
 & \text{if $n=4$} \\
N_k(n) - |U_{\text{out}}|-|U_{\text{in}}|-|P_{\text{out}}|-|P_{\text{in}}| \\
~~~~+~|U_{\text{out}}\cap P_{\text{in}}|+|P_{\text{out}}\cap U_{\text{in}}| \\
~~~~+~|U_{\text{out}}\cap U_{\text{in}}|+|P_{\text{out}}\cap P_{\text{in}}|,
 & \text{if $n>4$.}
\end{cases} \]
where $|X\cap Y|$ denotes the maximum possible cardinality for such a set, and $N_k(n)$ denotes the
number of non-negasymmetric $k$-ary $n$-tuples (see
Lemma~\ref{lemma:number_non-negasymmetric_tuples}).
\end{lemma}

\begin{theorem}  \label{theorem:new_NOS_bounds}
Suppose $k>2$ and $S=(s_i)$ is an $\mathcal{NOS}_k(n)$ of period $m$.
\begin{itemize}
\item[i)] If $n=2$ then:
\[ m \leq \begin{cases}
\frac{k^2-k}{2}, & \text{if $k$ is odd} \\
\frac{k^2-k-2}{2}, & \text{if $k$ is even}
\end{cases} \]
\item[ii)] If $n=3$ then:
\[ m \leq \begin{cases}
\frac{k^3-2k+1}{2}, & \text{if $k$ is odd} \\
\frac{k^3-2k-6}{2}, & \text{if $k$ is even}
\end{cases} \]
\item[iii)] If $n=4$ then:
\[ m \leq \begin{cases}
\frac{k^4-2k^2+1}{2}, & \text{if $k$ is odd} \\
\frac{k^4-2k^2+k-2}{2}, & \text{if $k$ is even}
\end{cases}\]
\item[iv)] If $n>4$ then:
\[ m \leq \begin{cases}
\frac{k^n-5k^{(n-1)/2}+4k}{2}, & \text{if $n$ and $k$ are odd} \\
\frac{k^n-6k^{(n-1)/2}+3k+2}{2}, & \text{if $n$ is odd and $k$ is even} \\
\frac{k^n-3k^{n/2}-2k^{(n-2)/2}+k^2+3k}{2}, & \text{if $n$ is even and $k$ is odd} \\
\frac{k^n-3k^{n/2}+k^2+k-2}{2}, & \text{if $n$ and $k$ are even.}
\end{cases}\]
\end{itemize}
\end{theorem}

\begin{proof}
The proof builds on, and uses the notation of, Lemma~\ref{lemma:NOS-bounds-high-level}.
\begin{itemize}

\item[i)] {\bf $n=2$}.  If $k$ is odd then $P_{\text{out}}=\emptyset$ by
    Corollary~\ref{corollary:excluded-tuples-by-parity}(i), since there are no negasymmetric
    non-uniform 1-tuples.  If $k$ is even then, by
    Corollary~\ref{corollary:excluded-tuples-by-parity}(iii), $|P_{\text{out}}|=2$, since there
    are 2 uniform-alternating negasymmetric 1-tuples by
    Lemma~\ref{lemma:NOS_tuple_numbers}(vi). The result follows from
    Lemma~\ref{lemma:NOS-bounds-high-level}.


\item[ii)] {\bf $n=3$}.  If $k$ is odd then, by
    Corollary~\ref{corollary:excluded-tuples-by-parity}(i), $|P_{\text{in}}|=|P_{\text{out}}|$
    equals the number of negasymmetric non-uniform 2-tuples, i.e.\ $k-1$ by
    Lemma~\ref{lemma:NOS_tuple_numbers}(v),(vii). In this case $|P_{\text{in}}\cap
    P_{\text{out}}|=k-1$ by Lemma~\ref{lemma:NOS-constraint-interactions}(i) and
    Lemma~\ref{lemma:NOS_tuple_numbers}(ii),(iv).

    If $k$ is even then, by Corollary~\ref{corollary:excluded-tuples-by-parity}(ii),
    $|P_{\text{in}}|=|P_{\text{out}}|$ equals the number of uniform or alternating
    negasymmetric 2-tuples where every element is either 0 or $k/2$, i.e.\ 4. Next observe that
    when evaluating $|P_{\text{in}}\cap P_{\text{out}}|$ we need only consider alternating
    2-tuples in which every element is either 0 or $k/2$, since an edge can only be outgoing
    from a uniform negasymmetric ($n-1)$-tuple and incoming to a uniform negasymmetric
    ($n-1)$-tuple if it corresponds to a uniform negasymmetric $n$-tuple, and such an edge
    cannot occur in $B^-_k(n-1)$.  Then $|P_{\text{in}}\cap P_{\text{out}}|$ is equal to the
    number of alternating negasymmetric $n$-tuples in which every element is either 0 or $k/2$,
    i.e.\ 2.

    The result follows from Lemma~\ref{lemma:NOS-bounds-high-level}.


\item[iii)] {\bf $n=4$}. By Corollary~\ref{corollary:NOS_excluded_tuples_by_unequal_degree}(v),
    $|U_{\text{out}}|$ equals the number of non-uniform-alternating left-sns $3$-tuples, i.e.\
    $k^2-k$ by Lemma~\ref{lemma:NOS_tuple_numbers}(x).  If $k$ is odd then, by
    Corollary~\ref{corollary:excluded-tuples-by-parity}(i), $|P_{\text{out}}|$ equals the
    number of negasymmetric non-uniform 3-tuples, i.e.\ $k-1$ by
    Lemma~\ref{lemma:NOS_tuple_numbers}(i),(v). If $k$ is even then, by
    Corollary~\ref{corollary:excluded-tuples-by-parity}(iii), $|P_{\text{out}}|$ equals the
    number of uniform-alternating negasymmetric 3-tuples, i.e.\ 2 by
    Lemma~\ref{lemma:NOS_tuple_numbers}(vi). The result follows from
    Lemma~\ref{lemma:NOS-bounds-high-level}.


\item[iv)a)] {\bf $n>4$; $n$ and $k$ odd}.  By
    Corollary~\ref{corollary:NOS_excluded_tuples_by_unequal_degree}(i),(ii),
    $|U_{\text{out}}|=|U_{\text{in}}|$ equals the number of non-uniform left-sns
    $(n-1)$-tuples, i.e.\ $k^{(n-1)/2}-1$ by Lemma~\ref{lemma:NOS_tuple_numbers}(ix).

    Also, by Lemma~\ref{lemma:NOS-constraint-interactions}(iv), $|U_{\text{out}}\cap
    U_{\text{in}}|=k-1$ since if $n\equiv 1\pmod 4$ there are $k$ choices for $c_0$, and one
    choice each for $c_1$ and $c_3$, giving $k$ possibilities, one of which is uniform (when
    $c_0=c_1=c_3=0$). An analogous argument applies if $n\equiv 3 \pmod 4$.

    By Corollary~\ref{corollary:excluded-tuples-by-parity} (i),
    $|P_{\text{out}}|=|P_{\text{in}}|$ equals the number of negasymmetric non-uniform
    $(n-1)$-tuples, i.e.\ $k^{(n-1)/2}-1$ by Lemma~\ref{lemma:NOS_tuple_numbers}(i),(v).

    By Lemma~\ref{lemma:NOS-constraint-interactions}(ii),(iii), $|U_{\text{out}}\cap
    P_{\text{in}}|=|P_{\text{out}}\cap U_{\text{in}}|=k-1$ (removing the case where the $c_i$
    are all equal).

    By Lemma~\ref{lemma:NOS-constraint-interactions}(i), $|P_{\text{out}}\cap P_{\text{in}}|$
    equals the number of negasymmetric non-uniform $(n-1)$-tuples that are also
    uniform-alternating, i.e.\ $k-1$ by Lemma~\ref{lemma:NOS_tuple_numbers}(iv),(vi).

    The result follows from Lemma~\ref{lemma:NOS-bounds-high-level}.


\item[iv)b)] {\bf $n>4$; $n$ odd and $k$ even}.  By
    Corollary~\ref{corollary:NOS_excluded_tuples_by_unequal_degree}(iii),(iv),
    $|U_{\text{out}}|=|U_{\text{in}}|$ equals the number of non-uniform non-alternating
    left-sns $(n-1)$-tuples, i.e.\ $2k^{(n-1)/2}-4$ by Lemma~\ref{lemma:NOS_tuple_numbers}(xi).

    Also, by Lemma~\ref{lemma:NOS-constraint-interactions}(iv), $|U_{\text{out}}\cap
    U_{\text{in}}|=4k-4$ since if $n\equiv 1\pmod 4$ there are $k$ choices for $c_0$, and two
    choices each for $c_1$ and $c_3$, giving $4k$ possibilities, two of which are uniform and
    two of which are alternating. An analogous argument applies if $n\equiv 3 \pmod 4$.

    By Corollary~\ref{corollary:excluded-tuples-by-parity}(ii),
    $|P_{\text{out}}|=|P_{\text{in}}|$ equals the number of uniform or alternating
    negasymmetric $(n-1)$-tuples, i.e.\ $k$ by Lemma~\ref{lemma:NOS_tuple_numbers}(v),(vii).

    By Corollary~\ref{corollary:excluded-tuples-by-parity}(ii), $P_{\text{out}}$ and
    $P_{\text{in}}$ only contain edges out-going/in-going from/to uniform or alternating
    $(n-1)$-tuples, and by
    Corollary~\ref{corollary:NOS_excluded_tuples_by_unequal_degree}(iii),(iv), $U_{\text{out}}$
    and $U_{\text{in}}$ only contain edges that are out-going/in-going from/to non-uniform
    non-alternating $(n-1)$-tuples, and hence $|U_{\text{out}}\cap
    P_{\text{in}}|=|P_{\text{out}}\cap U_{\text{in}}|=\emptyset$.

    By Corollary~\ref{corollary:excluded-tuples-by-parity}(ii), $P_{\text{out}}$ and
    $P_{\text{in}}$ only contain edges out-going/in-going from/to uniform or alternating
    negasymmetric $(n-1)$-tuples, which by Lemma~\ref{lemma:NOS_tuple_numbers}(vii) are
    uniform-alternating. Hence an edge in $P_{\text{out}}\cap P_{\text{in}}$ must be
    uniform-alternating.  An edge in $P_{\text{out}}\cap P_{\text{in}}$ cannot be
    negasymmetric, and there are $k-2$ non-negasymmetric uniform-alternating $n$-tuples by
    Lemma~\ref{lemma:NOS_tuple_numbers}(iii),(v).  Hence$|P_{\text{out}}\cap
    P_{\text{in}}|=k-2$.

    The result follows from Lemma~\ref{lemma:NOS-bounds-high-level}.

\item[iv)c)] {\bf $n>4$; $n$ even}. By
    Corollary~\ref{corollary:NOS_excluded_tuples_by_unequal_degree}(v),(vi),
    $|U_{\text{out}}|=|U_{\text{in}}|$ equals the number of non-uniform-alternating left-sns
    $(n-1)$-tuples, i.e.\ $k^{n/2}-k$ by Lemma~\ref{lemma:NOS_tuple_numbers}(x).

    Additionally, by Lemma~\ref{lemma:NOS-constraint-interactions}(iv), $|U_{\text{out}}\cap
    U_{\text{in}}|=k(k-1)$, by choosing $c_0$ and $c_2$ to be distinct.

    {\bf If $k$ is odd} then, by Corollary~\ref{corollary:excluded-tuples-by-parity}(i),
    $|P_{\text{out}}|=|P_{\text{in}}|$ equals the number of negasymmetric non-uniform
    $(n-1)$-tuples, i.e.\ $k^{(n-2)/2}-1$ by Lemma~\ref{lemma:NOS_tuple_numbers}(i),(v).
     \newline By Lemma~\ref{lemma:NOS-constraint-interactions}(ii),(iii), $|U_{\text{out}}\cap
    P_{\text{in}}|=|P_{\text{out}}\cap U_{\text{in}}|=k-1$, by ensuring that $c_0$,
    $c_1$ and $c_2$ are not all equal to $0$.
     \newline By Lemma~\ref{lemma:NOS-constraint-interactions}(i), $|P_{\text{out}}\cap
    P_{\text{in}}|=0$, i.e.\ the number of non-uniform alternating $(n-1)$-tuples in which
    every entry is either 0 or $k/2$ when $k$ is odd.

    {\bf If $k$ is even} then, by Corollary~\ref{corollary:excluded-tuples-by-parity}(iii),
    $|P_{\text{out}}|=|P_{\text{in}}|$ equals the number of uniform-alternating negasymmetric
    $(n-1)$-tuples, i.e.\ 2 by Lemma~\ref{lemma:NOS_tuple_numbers}(vi).
     \newline By Corollary~\ref{corollary:excluded-tuples-by-parity}(iii), $P_{\text{out}}$
    and $P_{\text{in}}$ respectively only contain edges out-going from or in-going to
    uniform-alternating $(n-1)$-tuples, and by
    Corollary~\ref{corollary:NOS_excluded_tuples_by_unequal_degree}(v),(vi), $U_{\text{out}}$
    and $U_{\text{in}}$ only contain edges that are out-going/in-going from/to
    non-uniform-alternating $(n-1)$-tuples, and hence $|U_{\text{out}}\cap
    P_{\text{in}}|=|P_{\text{out}}\cap U_{\text{in}}|=\emptyset$.
     \newline Finally, it is immediate that $|P_{\text{out}}\cap P_{\text{in}}|=2$,
    i.e.\ the number of non-uniform alternating $(n-1)$-tuples in which every entry is either 0
    or $k/2$ when $k$ is even.

    The result follows from Lemma~\ref{lemma:NOS-bounds-high-level}.

\end{itemize}
\end{proof}

\end{document}
